\section{A topology on ordinary singularity types}
\label{sec:ordinary-type-topology}

The two envelopes retain different parts of a singularity. Requiring both to converge makes their equality a closed condition. Controlled ideal inclusions and total variation of divisorial orders supply additional coordinates that the dyadic approximation can satisfy in a fixed Hermitian background. Together they give an analytic-closure criterion on the whole space of ordinary types.

Fix a Hermitian form $\omega$ on connected compact $X$ and write
\[
 \mathfrak S(X,\omega)=\PSH(X,\omega)/{\sim}.
\]
In this section $[u]$ denotes the ordinary bounded-difference class. For each smooth compact modification $\pi:Y\to X$, fix a Hermitian form $\gamma_\pi$ on $Y$. For a prime divisor $E\subset Y$, set $w_E=\int_E\gamma_\pi^{n-1}$ and use the generic number $\nu_E(u)=\nu(u\circ\pi,E)$ of \cref{sec:psh-divisorial-estimates}.

\begin{lemma}\label{lem:fixed-model-divisorial-variation}\index{Divisor!total variation on a fixed model}
For $u,v\in\PSH(X,\omega)$, the sum
\begin{equation}\label{eq:fixed-model-divisorial-variation}
 d_\pi([u],[v])=\sum_{E\subset Y\text{ prime}}w_E|\nu_E(u)-\nu_E(v)|
\end{equation}
is finite and defines a pseudometric on $\mathfrak S(X,\omega)$. Changing $\gamma_\pi$ gives an equivalent pseudometric on this fixed model.
\end{lemma}
\begin{proof}
The current $T=\pi^*\omega+\ddc(u\circ\pi)$ is positive. Divisorial removal in \cref{lem:nonclosed-divisorial-removal} gives, for every finite family of prime divisors,
\[
 \sum_Ew_E\nu_E(u)\leq\int_YT\wedge\gamma_\pi^{n-1}<\infty.
\]
Taking the supremum over finite families proves summability and countability of the positive coefficients. Apply the same bound to $v$. Bounded changes preserve the coefficients, and the triangle inequality holds term by term. On compact $Y$, two Hermitian forms bound one another by fixed positive multiples, so their divisor weights do also.
\end{proof}

For $t>0$ and $a>1$, let
\begin{equation}\label{eq:controlled-ideal-entourage}
 \mathcal U(t,a)=\{([u],[v]):\mathcal I(atu)\subseteq\mathcal I(tv),\quad
                         \mathcal I(atv)\subseteq\mathcal I(tu)\}.
\end{equation}
These are relations between ordinary classes, since a bounded change leaves each ideal unchanged. Define two further pseudometrics by
\[
 \begin{split}
 d_O([u],[v])&=\|P_\omega[u]-P_\omega[v]\|_{L^1(\mathrm dV)},\\
 d_I([u],[v])&=\|P_\omega[u]_{\mathcal I}-P_\omega[v]_{\mathcal I}\|_{L^1(\mathrm dV)}.
 \end{split}
\]

\begin{proposition}\label{prop:ordinary-type-uniformity}\index{Singularity type!uniform topology}
Finite intersections of the relations $d_O<r$, $d_I<r$, $d_\pi<r$ and $\mathcal U(t,a)$, with $r,t>0$ and $a>1$, form a base of entourages for a uniformity on $\mathfrak S(X,\omega)$. Denote its topology by $\tau_{\omega,0}$. A net $[u_j]$ converges to $[u]$ precisely when both envelopes converge in $L^1$, $d_\pi([u_j],[u])\to0$ for every fixed $\pi$, and $([u_j],[u])\in\mathcal U(t,a)$ eventually for every fixed $t,a$. Smooth positive choices of $\mathrm dV$ and of the $\gamma_\pi$ do not change the uniformity.
\end{proposition}
\begin{proof}
All generators are symmetric and contain the diagonal. Multiplier ideals decrease with the weight: locally subtract an upper bound from a potential to see this from integrability. For $b=\sqrt a$, put $\mathcal V=\mathcal U(t,b)\cap\mathcal U(bt,b)$. If $(x,y),(y,z)\in\mathcal V$, then
\[
 \mathcal I(atx)\subseteq\mathcal I(bty)\subseteq\mathcal I(tz),
\]
and the same argument with $x,z$ exchanged proves $\mathcal V\circ\mathcal V\subseteq\mathcal U(t,a)$. Halving a pseudometric radius gives the corresponding composition refinement. A finite intersection of these refinements proves the uniformity axiom. Its convergence criterion is exactly eventual membership in each generating entourage. Bounded ratios of smooth positive volumes and \cref{lem:fixed-model-divisorial-variation} prove independence of the auxiliary choices.
\end{proof}

Each entourage involves finitely many fixed models. The definition does not require a countable cofinal family of modifications or simultaneous uniform convergence over all models.

\subsection{Analytic closure}

Let $\mathcal A_\omega$ be the ordinary classes with analytic singularities in the rational-coefficient convention of \cref{sec:analytic-types}. The curvature correction in \cref{cor:fixed-hermitian-analytic-density} preserves this convention. The next theorem also retains the ideal weights affected by that correction.

\Needspace{14\baselineskip}
\begin{theorem}\label{thm:ordinary-topology-analytic-closure}\index{Analytic singularity!closure in the ordinary-type topology}
The closure and sequential closure of $\mathcal A_\omega$ in $\tau_{\omega,0}$ are both
\begin{equation}\label{eq:ordinary-topology-analytic-closure}
 \mathcal G_\omega=\{[u]:P_\omega[u]=P_\omega[u]_{\mathcal I}\}.
\end{equation}
For every $u\in\PSH(X,\omega)$, the dyadic sequence $b_j$ of \cref{cor:fixed-hermitian-analytic-density} satisfies
\[
 [b_j]\longrightarrow[u]\text{ in }\tau_{\omega,0}
 \quad\Longleftrightarrow\quad [u]\in\mathcal G_\omega.
\]
The approximants may be chosen with one positive rational analytic coefficient on all charts.
\end{theorem}
\begin{proof}
Use the notation there: $m_j=2^j$, $a_j$ has type $(1/m_j,\mathcal I(m_ju))$, $c_j=(1+\epsilon_j)^{-1}\uparrow1$ is rational, and $b_j=c_ja_j$. Its ordinary and ideal envelopes agree and decrease to $P_\omega[u]_{\mathcal I}$. Thus the two envelope coordinates converge to those of $u$ precisely when $[u]\in\mathcal G_\omega$.

Fix $t>0$, $a>1$, and choose $1<s<a$. For large $j$, $ac_j>s$ and $m_j\geq st/(s-1)$. By \cref{lem:bergman-controlled-ideals},
\[
 \mathcal I(atb_j)=\mathcal I(atc_ja_j)
 \subseteq\mathcal I(sta_j)\subseteq\mathcal I(tu).
\]
Since $u\preceq a_j\preceq b_j$, the reverse relation $\mathcal I(atu)\subseteq\mathcal I(tb_j)$ holds for every $j$. The scale $c_j$ belongs inside the first ideal weight.

For a divisor $E$ on a fixed $\pi$, the Bergman divisorial estimate gives
\[
 c_j\max\left\{\nu_E(u)-\frac{A(E)}{m_j},0\right\}
 \leq\nu_E(b_j)\leq c_j\nu_E(u)\leq\nu_E(u).
\]
Here $A(E)$ is the log discrepancy, computed over a coordinate neighborhood of its center. Each coefficient converges to $\nu_E(u)$. The summability in \cref{lem:fixed-model-divisorial-variation} therefore gives $d_\pi([b_j],[u])\to0$ by dominated convergence. This proves every remaining coordinate of convergence.

Finally analytic classes have equal envelopes by \cref{prop:analytic-envelope-type}. The map $[u]\mapsto(P_\omega[u],P_\omega[u]_{\mathcal I})$ is continuous into $L^1\times L^1$, whose diagonal is closed. Thus any limit of analytic nets lies in $\mathcal G_\omega$. The constructed sequence gives the converse and the sequential assertion.
\end{proof}

\subsection{Separation and mass}

The topology is defined before identifying classes with their ordinary envelopes. Those envelopes can coincide for distinct ordinary types.

\begin{proposition}\label{prop:ordinary-topology-nonseparation}\index{Singularity type!nonseparation}
Two classes have the same $\tau_{\omega,0}$ neighborhoods if and only if their ordinary envelopes coincide. If $\dim X>0$, this topology is not $T_0$ and cannot be induced by a metric on $\mathfrak S(X,\omega)$.
\end{proposition}
\begin{proof}
The $d_O$ coordinate distinguishes different ordinary envelopes. Conversely, equal ordinary envelopes have equal multiplier ideals at every weight by \cref{cor:ordinary-preserves-i}. Their ideal envelopes and all divisorial coefficients therefore agree, so every generating entourage has identical sections at the two classes.

To produce distinct such classes, choose a cutoff logarithmic pole $w\leq0$ supported in a coordinate ball, with $w$ plurisubharmonic near its pole and smooth on the cutoff annulus. A small positive multiple makes $w\in\PSH(X,\omega)$. Put $\chi(s)=1-\sqrt{1-s}$ for $s\leq0$ and $u=\chi(w)$. The inequalities $0<\chi'\leq1/2$ and $\chi''\geq0$, applied first to bounded truncations, give $u\in\PSH(X,\omega)$. Near the pole, $u$ has magnitude comparable to $\sqrt{-\log|z|^2}$, so it is unbounded and $\mathrm e^{-tu}$ is integrable for every $t>0$. All its multiplier ideals are trivial and $P_\omega[u]_{\mathcal I}=0$.

For every $0<\delta\leq1$, the function $\delta s-\chi(s)$ is bounded above on $(-\infty,0]$. Hence $\delta w\preceq u$ and $\delta w$ is an ordinary-envelope competitor. Letting $\delta\downarrow0$ off the pole gives $P_\omega[u]=0$. Thus $[u]\ne[0]$ have the same neighborhoods. In dimension zero the connected manifold is a point and the ordinary-type space is a singleton.
\end{proof}

The extended measures of \cref{thm:arbitrary-extended-ma} give two canonical mass labels
\[
 \int_X\mathrm{MA}_\omega(P_\omega[u]),\qquad
 \int_X\mathrm{MA}_\omega(P_\omega[u]_{\mathcal I}).
\]
In the present positive background, both are at most $\int_X\omega^n$. For the ideal envelope this follows from \cref{cor:ordinary-preserves-i,prop:fixed-point-finite}. For the ordinary envelope, normalize $u\leq0$ and put $p_j=P_\omega(u+j,0)$ and $p=P_\omega[u]$. Then $p_j\to p$ in capacity by \cref{lem:monotone-psh-capacity}. For $k>0$, set $\eta_k=\min\{1,(u+k)_+\}$. If $j>k$, the contact with $u+j$ lies outside $\{\eta_k>0\}$, so \cref{prop:rooftop-measure} gives $\eta_k\mathrm{MA}_\omega(p_j)\leq\eta_k\omega^n$. Test against a nonnegative continuous function and apply \cref{lem:weighted-extended-lsc} to obtain the same inequality for $p$. Letting $k\to\infty$ proves $\mathrm{MA}_\omega(p)\leq\omega^n$, since the remaining set $\{u=-\infty\}$ is pluripolar.

The original representative can nevertheless have infinite mass: \cref{prop:goodness-infinite-mass} gives such a $u$ with both envelopes equal to zero. The topology and its analytic-closure criterion include this ordinary class. A mass condition on the representative is therefore unnecessary, while infinite canonical mass does not arise for these two envelopes in a positive Hermitian background.

\subsection{Transport with an obstacle}

For a smooth real obstacle $F$, use the ordinary envelope $P_{\omega,F}[u]$ of \cref{def:singularity-envelope}, and write $P_{\omega,F}[u]_{\mathcal I}$ for the ideal envelope with $v\leq F$ in place of $v\leq0$. Its competitor family contains a translate of $u$. Replacing a competitor by its maximum with this translate shows that the ideal envelope can equivalently use $v\preceq_{\mathcal I}u$. The same notation applies to any smooth real background admitting $u$. Define $\tau_{\omega,F}$ by replacing the two envelope coordinates in \cref{prop:ordinary-type-uniformity} by these obstacle envelopes.

\Needspace{10\baselineskip}
\begin{proposition}\label{prop:ordinary-topology-background-transport}
Let $\chi\in C^\infty(X,\mathbb R)$ and $\omega'=\omega+\ddc\chi>0$. The map $[u]\mapsto[u-\chi]$ is a uniform isomorphism
\[
 (\mathfrak S(X,\omega),\tau_{\omega,F})
 \longrightarrow(\mathfrak S(X,\omega'),\tau_{\omega',F-\chi}).
\]
For $s>0$, $[u]\mapsto[su]$ is a uniform isomorphism from $\tau_{\omega,F}$ to $\tau_{s\omega,sF}$. Both assertions transport the ordinary and ideal envelopes exactly.
\end{proposition}
\begin{proof}
For either envelope operator $P$, its competitors give
\[
 P_{\omega',F-\chi}[u-\chi]=P_{\omega,F}[u]-\chi,
 \qquad P_{s\omega,sF}[su]=sP_{\omega,F}[u].
\]
Smooth additions leave all ideals and divisorial numbers unchanged. Scaling multiplies the divisorial and envelope distances by $s$ and reindexes the ideal weights by $t\mapsto st$. Each map and its inverse therefore transport the generating entourages.
\end{proof}

For two smooth obstacles, translation of competitors gives
\[
 P_{\omega,F}[u]+\min_X(G-F)
 \leq P_{\omega,G}[u]
 \leq P_{\omega,F}[u]+\max_X(G-F),
\]
with the same inequalities for ideal envelopes. The associated extended measures satisfy $\mathrm{MA}_{\omega'}(u-\chi)=\mathrm{MA}_\omega(u)$ and $\mathrm{MA}_{s\omega}(su)=s^n\mathrm{MA}_\omega(u)$ by the local construction in \cref{thm:arbitrary-extended-ma}, including at infinite total mass. Keeping the zero obstacle fixed in both backgrounds instead gives $P_{\omega',0}[u-\chi]=P_{\omega,\chi}[u]-\chi$. The bounded comparison of the obstacles does not give continuity of this change.

\begin{proposition}\label{prop:ordinary-topology-gauge-failure}\index{Singularity type!failure of smooth-gauge continuity}
Let $Y$ be a compact connected Riemann surface and $\theta>0$ a smooth form of mass one. If $F$ is a nonconstant smooth function with $\theta'=\theta+\ddc F>0$, the bijection
\[
 \Gamma_F:(\mathfrak S(Y,\theta),\tau_{\theta,0})
 \longrightarrow(\mathfrak S(Y,\theta'),\tau_{\theta',0}),
 \qquad [u]\longmapsto[u-F],
\]
and its inverse are both discontinuous. Equivalently, $\tau_{\theta,0}$ and $\tau_{\theta,F}$ are incomparable topologies on the same ordinary-type space.
\end{proposition}
\begin{proof}
The smooth function $\theta'/\theta$ has mean one and is not constant, since a harmonic function on $Y$ is constant. Choose $\min_Y(\theta'/\theta)<c<1$, so $\theta'\not\geq c\theta$.

Choose nonatomic polar probability measures $\mu_j\rightharpoonup\theta$. To construct them, start with finite atomic probabilities converging weakly to $\theta$ as in \cref{thm:curve-ideal-lone-completion}. Replace each atom by a probability supported on a polar Cantor set in an arbitrarily small coordinate disc about that atom, preserving its weight. The construction in \cref{ex:diffuse-polar-zero-model} supplies such nonatomic probabilities in every disc. Choose the discs with diameters tending uniformly to zero. Uniform continuity of continuous test functions preserves weak convergence, and a finite union of their polar carriers is polar.

Let $q_j$ be the normalized Green potentials with $\theta+\ddc q_j=\mu_j$ and $\sup_Yq_j=0$. Normalized compactness gives $L^1$ cluster limits, and every such limit has curvature $\theta$ and supremum zero. It is therefore zero, so $q_j\to0$ in $L^1$. Put $u_j=cq_j$. Then
\[
 \theta+\ddc u_j=(1-c)\theta+c\mu_j,\qquad
 u_j\leq P_\theta[u_j]\leq0.
\]
All point Lelong numbers vanish. Since modifications of a smooth curve are isomorphisms, \cref{thm:local-i-inputs} gives trivial multiplier ideals at every weight. Thus $P_\theta[u_j]_{\mathcal I}=0$, all controlled ideal relations hold, and every divisorial distance to zero vanishes. The displayed bounds give ordinary-envelope convergence in $L^1$. Consequently $[u_j]\to[0]$ in $\tau_{\theta,0}$.

Set $v_j=P_{\theta,F}[u_j]$. For every competitor $h\preceq u_j$, \cref{cor:polar-order-maximum} gives $\theta+\ddc h\geq c\mu_j$. The family below $F$ contains a translate of $u_j$ and is stable under finite maxima. Choquet selection and distributional convergence of its increasing maxima retain this inequality at $v_j$. If $[u_j]\to[0]$ in $\tau_{\theta,F}$, its ordinary-envelope coordinate would give $v_j\to P_{\theta,F}[0]=F$ in $L^1$, because $F$ is a bounded $\theta$-plurisubharmonic competitor. Passing the current inequality to the limit would give $\theta'\geq c\theta$, a contradiction.

The exact paired transport in \cref{prop:ordinary-topology-background-transport} identifies this failure with discontinuity of $\Gamma_F$. Apply the same argument to $\theta'$ and $-F$ to prove discontinuity of its inverse. Under paired transport these are precisely the two directions of the identity between $\tau_{\theta,0}$ and $\tau_{\theta,F}$, proving incomparability.
\end{proof}

The sequence has trivial ideals and polar curvature $c\mu_j$. By \cref{lem:curve-polar-model} and the bounded obstacle comparison, its ordinary envelopes in both backgrounds have non-pluripolar mass $1-c$. The discontinuity therefore persists inside a positive-mass cutoff. On the same curve, the complete polar-variation metric of \cref{thm:curve-model-volume-distance} is invariant under smooth background changes after ordinary projection; it retains the diffuse curvature that the sequence above loses.

\subsection{All smooth obstacles}
\label{subsec:all-obstacle-topology}

The preceding failure comes from choosing one obstacle. Requiring convergence for every smooth obstacle makes the collection stable under a change of background. The same dyadic approximation still determines its analytic closure, now by testing envelope equality at every obstacle.

Let $\tau^{\mathrm{all}}_\omega$ be the uniform topology generated by all the uniformities $\tau_{\omega,F}$, with $F\in C^\infty(X,\mathbb R)$. Thus an entourage uses finitely many obstacle-envelope distances, together with finitely many of the ideal relations and divisorial distances already defined.

\begin{proposition}\label{prop:all-obstacle-uniformity}\index{Singularity type!all smooth obstacles}
A countable uniformly dense family $\{F_k\}_{k\geq1}\subset C^\infty(X,\mathbb R)$ suffices for the obstacle coordinates of $\tau^{\mathrm{all}}_\omega$. For $\chi$ smooth with $\omega'=\omega+\ddc\chi>0$, the map $[u]\mapsto[u-\chi]$ is a uniform isomorphism between $\tau^{\mathrm{all}}_\omega$ and $\tau^{\mathrm{all}}_{\omega'}$. Positive scaling gives the corresponding uniform isomorphism to $\tau^{\mathrm{all}}_{s\omega}$. In positive dimension this topology is not $T_0$ on ordinary classes.
\end{proposition}
\begin{proof}
For either envelope operator, the obstacle comparison gives
\[
 \|P_{\omega,F}[u]-P_{\omega,G}[u]\|_{L^1(\mathrm dV)}
 \leq\|F-G\|_\infty.
\]
The difference between the two envelope distances for a pair $[u],[v]$ is therefore at most $2\|F-G\|_\infty$, uniformly in the pair. Smooth functions contain a countable dense subset for the uniform norm: approximate each member of a countable dense subset of $C(X)$ by smooth functions with errors $1/m$, $m\geq1$. Approximating each of the finitely many obstacles in an entourage and then shrinking its radii proves equality of the uniformities. The divisorial models are still allowed to range over every fixed compact modification.

The exact transport in \cref{prop:ordinary-topology-background-transport} reindexes the entire smooth-obstacle family by $F\mapsto F-\chi$, or by $F\mapsto sF$. These are bijections of that family, proving both isomorphisms.

For nonseparation, use the unbounded $u=1-\sqrt{1-w}$ in \cref{prop:ordinary-topology-nonseparation}. It has trivial ideals and zero divisorial numbers, and $\delta w\preceq u$ for $0<\delta\leq1$. Fix $F$ and put $q=P_\omega(F)$. Since $\omega>0$, the constant $\min_XF$ is a competitor, so $q$ is bounded and $q\leq F$. The function
\[
 h_\delta=(1-\delta)q+\delta w-\delta\|F\|_\infty
\]
is $\omega$-plurisubharmonic, lies below $F$, and satisfies $h_\delta\preceq u$. Letting $\delta\downarrow0$ off the pole of $w$ gives $P_{\omega,F}[u]=q=P_{\omega,F}[0]$. Trivial ideals give the same equality for the ideal envelopes. Every generating entourage therefore has identical sections at $[u]$ and $[0]$, although these ordinary classes differ.
\end{proof}

\begin{theorem}\label{thm:all-obstacle-analytic-closure}\index{Analytic singularity!closure for all obstacles}
The closure and sequential closure of $\mathcal A_\omega$ in $\tau^{\mathrm{all}}_\omega$ both equal
\[
 \mathcal G^{\mathrm{all}}_\omega
 =\{[u]:P_{\omega,F}[u]=P_{\omega,F}[u]_{\mathcal I}
                      \text{ for every smooth }F\}.
\]
It suffices to test equality on the countable dense family above. The dyadic sequence $b_j$ of \cref{cor:fixed-hermitian-analytic-density} converges to $[u]$ in this topology if and only if $[u]\in\mathcal G^{\mathrm{all}}_\omega$.
\end{theorem}
\begin{proof}
Fix a smooth obstacle $F$. The increasing-maxima proof of \cref{thm:hermitian-i-projection} applies to the family below $F$: a translate of $u$ is a lower competitor, and the fixed upper bound supplies the local convergence needed for increasing-limit stability of ideals. Thus $P_{\omega,F}[u]_{\mathcal I}\sim_{\mathcal I}u$. For an analytic $a$, \cref{lem:analytic-ideal-dominance} makes this ideal envelope an ordinary competitor, so $P_{\omega,F}[a]=P_{\omega,F}[a]_{\mathcal I}$.

Write $p=P_{\omega,F}[u]_{\mathcal I}$ and $q_j=P_{\omega,F}[b_j]_{\mathcal I}$. The dyadic construction gives $u\preceq_{\mathcal I}b_j$, decreasing ordinary types, and convergence of every local divisorial Lelong number to that of $u$. Ideal order, proved by taking maxima with a translate of $b_j$ below $F$, gives $p\leq q_j\leq F$. Analytic envelope equality and the decreasing types make $q_j$ decrease. Its limit $q$ is genuine and $\omega$-plurisubharmonic, with $p\leq q\leq F$; the same bounds give convergence in $L^1$.

Pullback convergence and Lelong semicontinuity, as in \cref{prop:ordered-ideal-envelope-convergence}, give $\nu_E(q)\geq\nu_E(u)$ for every local divisor $E$. The inequality $p\leq q$ gives the reverse bound. The local ideal criterion yields $q\sim_{\mathcal I}u$. Since $q\leq F$, it is admissible for $p$, proving
\[
 P_{\omega,F}[b_j]=P_{\omega,F}[b_j]_{\mathcal I}
       \searrow P_{\omega,F}[u]_{\mathcal I}.
\]
The controlled ideal relations and all fixed-model divisorial distances converge by the proof of \cref{thm:ordinary-topology-analytic-closure}; they do not depend on $F$. This one dyadic sequence consequently converges in $\tau^{\mathrm{all}}_\omega$ exactly under the stated envelope equalities.

For each $F$, the equality of the two $L^1$ envelope coordinates is closed. Analytic classes belong to every such equality set, so their net closure lies in the displayed intersection. The constructed sequence proves the converse. Finally the uniform obstacle bound in the preceding proof extends equality from a dense family to every smooth obstacle.
\end{proof}

The inclusion $\mathcal G^{\mathrm{all}}_\omega\subseteq\mathcal G_\omega$ follows by taking $F=0$. On curves the two sets coincide, as proved in \cref{prop:curve-all-obstacle-topology}. The theorem does not identify them in higher dimension. Although $\tau^{\mathrm{all}}_\omega$ removes the fixed-obstacle gauge dependence, it still has incomplete curve families and remains strictly weaker than polar total variation on ordinary curve models, by \cref{prop:curve-all-obstacle-incomplete}.

